\chapter{Backend}
\label{chap:backend}

\section{Monolithe modulaire}

Le backend (\repo{logiflow-backend}) est un monolithe \textbf{Spring Modulith} : un seul processus déployable, plusieurs modules métier aux frontières vérifiées en CI (ArchUnit + tests Modulith). Les modules ne se parlent que via leur package \texttt{api} ; \texttt{shared} est le seul module ouvert.

\figmmd{modules-backend}{Graphe de dépendances entre modules Spring Modulith (flèches vers \texttt{shared} omises).}

\begin{table}[ht]
\centering
\caption{Modules métier et responsabilité.}
\label{tab:modules}
\begin{tabular}{@{}lp{8.2cm}@{}}
\toprule
Module & Responsabilité \\
\midrule
\texttt{shared} & VOs, erreurs RFC~7807, audit, helpers sécurité \\
\texttt{iam} & Utilisateurs et rôles applicatifs \\
\texttt{referential} & Sites, clients, marchandises \\
\texttt{fleet} & Véhicules et remorques \\
\texttt{driver} & Chauffeurs, conformité, disponibilités \\
\texttt{order} & Commandes et lignes \\
\texttt{dossier} & Dossiers de transport, segments \\
\texttt{planning} & Voyages, trajets, affectations, conformité \\
\texttt{tracking} & Événements d'exécution \\
\texttt{maintenance} & OT, plans, sinistres, scores de santé \\
\texttt{carburant} & Prises, stations \\
\texttt{document} & Pièces jointes polymorphes \\
\texttt{ai} & Façade BFF vers Flask, journal des interactions \\
\texttt{config} & Security, CORS, async, scheduling \\
\bottomrule
\end{tabular}
\end{table}

\section{Architecture hexagonale par module}

Chaque module suit le même découpage : \texttt{domain} (agrégats, règles), \texttt{application} (cas d'usage), \texttt{infrastructure} (JPA, REST, adapters), \texttt{api} (contrat sortant vers les autres modules).

\figmmd{hexagonal}{Architecture hexagonale typique d'un module LogiFlow.}

\section{Modèle de données logique}

Les schémas SQL sont séparés par module (\texttt{commande}, \texttt{dossier}, \texttt{planning}, \ldots). Les liens inter-modules sont des UUID / JSON, pas des \texttt{@ManyToOne} JPA transverses.

\figmmd{er-logique}{Schéma entité-relation logique du cœur TMS et des satellites flotte / carburant / maintenance.}

\section{API et erreurs}

L'API publique est versionnée sous \texttt{/api/v1/**}. Les problèmes métier et techniques sont renvoyés au format \textbf{RFC~7807} (\texttt{application/problem+json}) \cite{rfc7807}. Les migrations de schéma passent par \textbf{Flyway}.

\section{Événements inter-modules}

Les publications utilisent \texttt{ApplicationEventPublisher} et \texttt{@ApplicationModuleListener} (asynchrone, après commit). Exemples : confirmation de commande $\rightarrow$ ouverture de dossier ; immobilisation atelier $\rightarrow$ indisponibilité flotte pour le planning.
